\section{Introduction}

A Riemannian metric has \emph{no conjugate points} if no nonzero Jacobi
field along a geodesic vanishes at two distinct times. Nonpositive
sectional curvature implies this condition. Gulliver showed that the
converse fails for an individual metric: local modifications of negatively
curved metrics can produce regions of positive curvature while preserving
the absence of conjugate points \cite[Section~3]{Gulliver1975}. Those
examples retain a manifold that already admits a negatively curved
metric. The existence question is therefore different: must a closed
manifold carrying a metric without conjugate points also carry some
metric of nonpositive sectional curvature?

This question appears in Ivanov and Kapovitch
\cite[Questions~1.1 and~8.1]{IvanovKapovitch2014} and in the survey of
Burns and Matveev \cite[Question~4.1.2]{BurnsMatveev2021}, which explicitly
includes dimension three. A closed surface without conjugate points
has universal cover diffeomorphic to $\mathbb R^2$ by the exponential
map, and hence has nonpositive Euler characteristic. Uniformization
then supplies a flat or negatively curved metric, so the existence
problem begins in dimension three. Ivanov and Kapovitch identify graph manifolds
as a central three-dimensional test case and ask about gluings of two
punctured-torus products \cite[Section~8]{IvanovKapovitch2014}.
We give a negative answer in dimension three, using a metric on a
two-piece graph manifold whose topology obstructs nonpositive curvature.

\begin{theorem}\label{thm:main}
There exists a closed connected orientable smooth three-manifold $M$ with a $C^\infty$ Riemannian metric without conjugate points, such that $M$ admits no $C^\infty$ Riemannian metric whose sectional curvature is everywhere at most zero.
\end{theorem}

The obstruction is stronger than the Riemannian assertion in the theorem.
The group $\pi_1(M)$ admits no proper cocompact isometric action on a
nonempty proper $\operatorname{CAT}(0)$ space
(Theorem~\ref{thm:no-cat0-action}); consequently, $M$ admits no locally
$\operatorname{CAT}(0)$ length metric inducing its topology
(Corollary~\ref{cor:no-local-cat0}).

The same manifold also admits no metric without focal points.
Recall that a metric has \emph{no focal points} if, for every nonzero
normal Jacobi field $J$ along a unit-speed geodesic with $J(0)=0$,
one has $\frac{d}{ds}|J(s)|^2>0$ for $s>0$. This condition is stronger
than absence of conjugate points. Ivanov and Kapovitch prove that a
closed three-manifold admits a metric without focal points if and only
if it admits a metric of nonpositive sectional curvature
\cite[Theorem~7.1]{IvanovKapovitch2014}. Their equivalence and
Theorem~\ref{thm:main} give the stated additional obstruction.

The manifold is obtained by gluing two copies of $\Sigma\times S^1$,
where $\Sigma$ is a compact oriented surface of genus one with one
boundary component. Write $h_i$ for the surface boundary and $f_i$ for
the circle factor in the $i$th boundary torus. The gluing is
\[
 h_2=h_1+f_1,\qquad f_2=h_1.
\]
This is an instance of Leeb's two-piece graph-manifold obstruction
\cite[Section~4.1, Example~4.1]{Leeb1995}. We include the group-theoretic
proof, which also gives the stronger action obstruction above.
Under a hypothetical nonpositively curved metric, the common torus
group acts by translations on a Euclidean plane. In each product piece,
the surface boundary is a commutator of elements commuting with the
fiber. The resulting translation vectors must be perpendicular.
These two orthogonality conditions are incompatible with the gluing.

The main work is the metric construction and its geodesic estimate.
Each product piece has nonpositive sectional curvature. Between them
we insert a neck $[0,1]\times T^2$ with metric
\[
 g=dt^2+dx^{\mathsf T}P(t)^{-1}\,dx,
\]
where $P(t)$ is a smooth path of positive definite matrices. We call
$P$ the \emph{cometric}, since it is the inverse torus metric. We choose
it concave, meaning that $-P''$ is positive semidefinite, with exact
product formulas near the ends. A small parameter $\delta$ controls
its derivatives; a shear of size $\delta^2$ realizes the integral
gluing after the torus periods are chosen. Convex warped collars then
join the endpoints smoothly to punctured-torus cores.

For the absence of conjugate points, we use the classical index-form
framework \cite[Sections~13--15]{Milnor1963}. Concavity of $P$ makes
the coordinate second variation nonnegative, but the intrinsic index
form differs by endpoint terms. The estimates must allow arbitrary
field values at the ends of an entire neck passage, so that they can
be combined with the estimates in the product pieces. A tangential
change of field handles passages with no turn. A returning geodesic
requires an estimate over its whole excursion.

For such an excursion, a positive boundary term controls one component
of the endpoint change in the torus part of the variation field.
We combine this control with the coordinate energy by Hilbert-space
Cauchy--Schwarz and a quadratic minimization. A strict timing margin
for excursions leaving an endpoint collar makes the estimate uniform
in the entry angle, turning height, and excursion time. No upper
bound on the duration is needed. This is what permits one metric to
handle every complete geodesic.

Section~\ref{sec:topology} fixes the manifold and proves its curvature
obstruction. Section~\ref{sec:cometric} constructs the metric and
establishes the matrix estimates. Section~\ref{sec:variations} derives
the index identities and classifies neck passages, and
Section~\ref{sec:turning} proves the returning-passage estimate.
Section~\ref{sec:global} sums the intrinsic index integrals of the
original field over all visits, including infinitely many visits,
and completes the proof of Theorem~\ref{thm:main}.
Appendix~\ref{app:punctured-torus} records the hyperbolic cusp model
used for the compact product cores.
